\section{Summary}\label{summary}

PyTorch (Paszke et al. 2019) is the dominant framework for machine
learning research and deployment, providing tensor operations and
automatic differentiation used across virtually all modern deep learning
workflows. However, PyTorch's high-performance execution has
historically been tied to specific hardware ecosystems --- primarily
NVIDIA GPUs via CUDA and Apple Silicon via MPS --- requiring separate
implementations for each target platform.

WebGPU (W3C GPU for the Web Working Group 2024) is a modern graphics and
compute API that provides a unified abstraction over diverse GPU
hardware, with implementations available across desktop, mobile, and
browser environments. By targeting WebGPU, computations become portable
across hardware that would otherwise require distinct backends.

torch-webgpu bridges these two ecosystems. It implements a PyTorch
out-of-tree backend that compiles PyTorch models to WebGPU compute
operations, enabling PyTorch tensors and models to run on any
WebGPU-capable hardware. Tensor operations are implemented as WGSL
(WebGPU Shading Language) compute shaders, and the backend integrates
with PyTorch's \texttt{torch.compile} infrastructure via the
\texttt{PrivateUse1} device extension mechanism introduced in PyTorch
2.0.

\section{Statement of need}\label{statement-of-need}

High-performance machine learning research depends on PyTorch, but
hardware portability remains a significant challenge. Running PyTorch
models on hardware without CUDA support --- including integrated GPUs,
mobile devices, and non-NVIDIA discrete GPUs --- currently requires
either CPU fallback with substantial performance loss, or substantial
engineering effort to port to a new backend.

WebGPU provides a hardware abstraction layer with broad support across
GPU vendors and platforms, including browser-based execution. However,
no existing tool provides a direct path from PyTorch models to WebGPU
execution. Solutions such as TensorFlow.js (Abadi et al. 2016) and ONNX
Runtime Web operate on different model representations and require
format conversion, breaking the native PyTorch workflow.

torch-webgpu fills this gap by implementing WebGPU as a native PyTorch
device. Researchers can run existing PyTorch models on WebGPU hardware
using the standard \texttt{torch.compile} interface and
\texttt{device="webgpu"} semantics, without modifying model code or
converting to intermediate formats.

\section{State of the field}\label{state-of-the-field}

PyTorch 2.0 introduced \texttt{torch.compile} with TorchDynamo and
TorchInductor as its default compiler stack, targeting CUDA and CPU. The
\texttt{PrivateUse1} mechanism allows out-of-tree backends to register
as first-class PyTorch devices. Existing out-of-tree backends include
implementations for Ascend NPUs (Huawei Ascend 2024) and other
vendor-specific hardware, demonstrating the viability of this approach.

On the WebGPU side, TensorFlow.js provides GPU-accelerated ML in
browsers using WebGL and WebGPU, but operates independently of the
PyTorch ecosystem. ONNX Runtime Web supports inference from ONNX models
in browsers, but again requires format conversion from PyTorch. Neither
provides the native PyTorch device experience that torch-webgpu targets.

\section{Software design}\label{software-design}

torch-webgpu is implemented primarily in C++ with a thin Python wrapper
distributed as a PyPI package. The backend is structured around four
components.

The WebGPU context (\texttt{webgpu\_context.cpp}) manages adapter and
device lifecycle using Dawn (Google 2024), Google's native WebGPU
implementation. On initialization, the context requests a
high-performance adapter and falls back to any available adapter if none
is found, then queries adapter limits to configure the device with
maximum supported capabilities.

The WebGPU allocator (\texttt{webgpu\_allocator.cpp}) implements
PyTorch's allocator interface, mapping tensor memory to WebGPU buffer
objects with \texttt{Storage\ \textbar{}\ CopySrc\ \textbar{}\ CopyDst}
usage flags. Buffer sizes are aligned to 4-byte boundaries as required
by WebGPU's \texttt{WriteBuffer} operation.

PyTorch operator bindings (\texttt{bindings.cpp}) register the backend
under PyTorch's \texttt{PrivateUse1} device type using
\texttt{TORCH\_LIBRARY\_IMPL}. Operators with native WebGPU
implementations are registered directly; unsupported operators fall back
to CPU via \texttt{at::native::cpu\_fallback}.

Compute operations are implemented as WGSL compute shaders. Activation
functions including ReLU, GeLU, and SiLU
(\texttt{activation\_functions.cpp}) are dispatched through PyTorch's
\texttt{TensorIterator} infrastructure to WGSL kernels. Matrix
multiplication (\texttt{mm.wgsl}) uses a 2D dispatch where each GPU
thread computes one output element of C = A x B. The shader accepts a
\texttt{Params} uniform buffer containing matrix dimensions, per-tensor
storage offsets, and stride arrays of up to 8 dimensions, enabling
correct indexing of non-contiguous and transposed tensors without
requiring data copies. Workgroup size is fixed at 16x16 threads, with
dispatch dimensions computed to cover the full M x K output space.

\section{Research impact statement}\label{research-impact-statement}

torch-webgpu has been used to conduct empirical benchmarking of PyTorch
model inference across heterogeneous hardware via WebGPU, the results of
which are reported in (Maczan 2026). That study demonstrates that
WebGPU-based execution enables cross-platform ML inference on hardware
not supported by existing PyTorch backends.

\section{AI usage disclosure}\label{ai-usage-disclosure}

Anthropic's Claude Opus 4.6 was used to generate and optimize some WGSL
shader code and to generate test scripts. All design decisions,
architecture, and research were made by the author. Claude was also used
to proof-read and improve the quality of this paper. All AI-assisted
outputs were reviewed and validated by the author.

\section{Acknowledgements}\label{acknowledgements}

The author thanks the PyTorch team for their documentation on
out-of-tree backends; Elie Michel for the WebGPU C++ guide (Michel
2024); the W3C WebGPU working group for the WGSL specification; the
WebGPU Fundamentals project (WebGPU Fundamentals Contributors 2024) for
educational resources; and the Ascend PyTorch team (Huawei Ascend 2024)
for a high-quality reference implementation of a PrivateUse1 backend.

\section*{References}\label{references}
\addcontentsline{toc}{section}{References}

\phantomsection\label{refs}
\begin{CSLReferences}{1}{0}
\bibitem[\citeproctext]{ref-abadi2016tensorflow}
Abadi, Martı́n, Paul Barham, Jianmin Chen, Zhifeng Chen, Andy Davis,
Jeffrey Dean, Matthieu Devin, et al. 2016. {``{TensorFlow}: A System for
{Large-Scale} Machine Learning.''} In \emph{12th USENIX Symposium on
Operating Systems Design and Implementation (OSDI 16)}, 265--83. USENIX
Association.

\bibitem[\citeproctext]{ref-dawn2024}
Google. 2024. {``Dawn: A WebGPU Implementation.''} 2024.
\url{https://dawn.googlesource.com/dawn}.

\bibitem[\citeproctext]{ref-ascend2024}
Huawei Ascend. 2024. {``PyTorch Ascend: Out-of-Tree PrivateUse1 Backend
for Ascend NPUs.''} 2024. \url{https://github.com/ascend/pytorch}.

\bibitem[\citeproctext]{ref-maczan2026webgpu}
Maczan, Jędrzej. 2026. {``Characterizing {WebGPU} Dispatch Overhead for
{LLM} Inference Across Four {GPU} Vendors, Three Backends, and Three
Browsers.''} arXiv preprint arXiv:2604.02344.
\url{https://arxiv.org/abs/2604.02344}.

\bibitem[\citeproctext]{ref-michel2024}
Michel, Elie. 2024. {``Learn WebGPU: A c++ Guide to WebGPU for Native
Graphics.''} 2024.
\url{https://eliemichel.github.io/LearnWebGPU/index.html}.

\bibitem[\citeproctext]{ref-paszke2019pytorch}
Paszke, Adam, Sam Gross, Francisco Massa, Adam Lerer, James Bradbury,
Gregory Chanan, Trevor Killeen, et al. 2019. {``PyTorch: An Imperative
Style, High-Performance Deep Learning Library.''} In \emph{Advances in
Neural Information Processing Systems}, 32:8024--35. Curran Associates,
Inc.

\bibitem[\citeproctext]{ref-webgpu2024}
W3C GPU for the Web Working Group. 2024. {``WebGPU Specification.''}
2024. \url{https://www.w3.org/TR/webgpu/}.

\bibitem[\citeproctext]{ref-webgpufundamentals2024}
WebGPU Fundamentals Contributors. 2024. {``WebGPU Fundamentals.''} 2024.
\url{https://webgpufundamentals.org/}.

\end{CSLReferences}
